\section{Internalizing the formula (drill it until it's yours)}

Understanding a formula once and \emph{owning} it are different. This section is
built for ownership: one memorable mental hook, the formula seen from four angles,
three worked drills of rising subtlety, and a blank-page recall test. Work through
it actively --- cover the answers, try first, then check.

\subsection{The one hook to remember everything}

If you remember \emph{one} sentence, remember this:

\begin{keybox}[The hook]
\centering\Large
\textbf{``Everyone shops the market \emph{minus themselves}; \\ add up what they
grab; never beat the smaller pile.''}
\end{keybox}

Every symbol in the formula is one clause of that sentence:
\[
M=\min\!\Big(\underbrace{\sum_{c}}_{\text{add up}}\ \underbrace{\min\big(b_c,}_{\text{what they'd grab}}\ \underbrace{S-s_c\big)}_{\text{market minus themselves}},\ \ \underbrace{\min(B,S)}_{\text{never beat the smaller pile}}\Big)
\]

Say the sentence out loud three times. Now the formula is just its transcription.

\subsection{The four angles (same truth, four viewpoints)}

Deep memory comes from seeing one thing many ways. Here is the formula as:

\begin{center}
\renewcommand{\arraystretch}{1.5}\footnotesize
\begin{tabular}{@{}p{2.4cm}p{11.4cm}@{}}
\toprule
\textbf{Angle} & \textbf{The formula in that language} \\
\midrule
\textbf{Words} & each firm buys from everyone but itself; total those; cap by the smaller side. \\
\textbf{Symbols} & $M=\min\big(\sum_c\min(b_c,\,S-s_c),\ \min(B,S)\big)$. \\
\textbf{Picture} & buyers on the left shake hands with \emph{other} firms' sellers on the right; count the hands. \\
\textbf{Code} & \code{reach=0; for c: reach += min(c.buy, totalSell - c.sell); return min(reach, min(B,S));} \\
\bottomrule
\end{tabular}
\end{center}

\begin{intu}
When one angle slips in the interview, another catches you. Forget the symbols?
Recite the sentence. Forget the sentence? Draw two columns and connect them.
Forget both? Write the four-line loop --- it \emph{is} the formula.
\end{intu}

\subsection{Why each piece \emph{must} be there (delete-it test)}

You truly own a formula when you know what breaks if you remove a piece.

\begin{center}
\renewcommand{\arraystretch}{1.35}\footnotesize
\begin{tabular}{@{}p{3.3cm}p{4.2cm}p{6.0cm}@{}}
\toprule
\textbf{Piece} & \textbf{If you drop it} & \textbf{The bug} \\
\midrule
the $-\,s_c$ & $\min(b_c,S)$ & a firm trades with \emph{itself}; self-only case wrongly $>0$ \\
the $\sum_c$ (per company) & $\min(B,S)$ & the naive answer; over-counts when one firm dominates \\
the outer $\min(B,S)$ & just $\sum_c\min(\dots)$ & overlapping reaches claim more trades than a side has \\
the inner $\min$ & $\sum_c b_c=B$ & a firm ``buys'' more than the market can sell it \\
\bottomrule
\end{tabular}
\end{center}

\begin{intu}
Each term is load-bearing. $-s_c$ enforces \emph{different company}. The
per-company sum enforces \emph{fairness across firms}. The two $\min$s enforce
\emph{you can't exceed what's available} --- once locally (per firm), once globally
(the smaller pile).
\end{intu}

\subsection{Drill 1 --- warm-up (cover the answer, compute first)}

\[
\Buy{}50(A),\qquad \Sell{}50(B).
\]
\emph{Your turn: totals? each firm's reach? answer?}

\smallskip
\noindent\textbf{Answer.} $B{=}50$, $S{=}50$. A wants $50$, allowed $S-s_A=50-0=50$,
reach $=50$. B has no buys. Sum $=50$; cap $\min(50,50)=50$. \textbf{$M=50$.} A
clean cross-company trade --- nothing blocks it.

\subsection{Drill 2 --- self-block plus a partial cross}

\[
\Buy{}80(A),\qquad \Sell{}80(A),\qquad \Sell{}30(B).
\]
\emph{Your turn before reading on.}

\smallskip
\noindent\textbf{Answer.} $B{=}80$, $S{=}80+30{=}110$.
\begin{itemize}[leftmargin=1.4em,nosep]
  \item A wants $80$; allowed $=S-s_A=110-80=30$; reach $=\min(80,30)=\mathbf{30}$.
  \item B has no buys; reach $0$.
\end{itemize}
Sum $=30$; cap $\min(80,110)=80$. \textbf{$M=30$.} A's own $80$ sells are off-limits
to it, so A can only reach B's $30$. \emph{This is the whole problem in one line:
the $-s_A$ chopped A's giant sell pile out of its own reach.}

\subsection{Drill 3 --- three firms, subtle}

\[
\Buy{}200(A),\ \Sell{}150(A),\ \Buy{}40(B),\ \Sell{}90(C).
\]
\emph{Try it. Tabulate first, then reach per firm, then cap.}

\smallskip
\noindent\textbf{Answer.} Totals: $B=200+40=240$, $S=150+90=240$.
\begin{center}\footnotesize
\begin{tabular}{@{}lllll@{}}
\toprule
firm & wants $b_c$ & own sells $s_c$ & allowed $S-s_c$ & reach $\min$ \\
\midrule
A & $200$ & $150$ & $240-150=90$ & $\mathbf{90}$ \\
B & $40$  & $0$   & $240-0=240$  & $\mathbf{40}$ \\
C & $0$   & $90$  & $240-90=150$ & $0$ \\
\bottomrule
\end{tabular}
\end{center}
Sum of reaches $=90+40+0=\mathbf{130}$. Cap $=\min(240,240)=240$, which does
\emph{not} bite. \textbf{$M=130$.}

\begin{intu}
Note what happened: total buy equals total sell ($240=240$), so a naive reader
would guess $240$. The truth is $130$, because A's big $200$ buy is throttled ---
it may only reach the $90$ of sells that aren't A's own. \emph{The reach term bound
the answer, not the cap.} If you can explain \emph{why 130 and not 240}, you own
the formula.
\end{intu}

\subsection{Blank-page recall test}

Close the document. On a blank page, in under a minute, produce:
\begin{enumerate}[leftmargin=1.6em,nosep]
  \item the one-sentence hook;
  \item the formula in symbols;
  \item the four-line code loop;
  \item why the $-s_c$ is there (one sentence);
  \item Drill 3's answer and \emph{why it's 130 not 240}.
\end{enumerate}
If all five come out, the formula is internalized. If one stalls, that's exactly
the piece to re-read --- then blank-page it again tomorrow. Spacing the recall by a
day is what moves it from short-term to permanent.

\subsection{The 20-second whiteboard routine}

Under interview pressure, don't recall --- \emph{reconstruct}. Always in this order:

\begin{enumerate}[leftmargin=1.6em,nosep]
  \item Draw the per-company table: firm, $b_c$, $s_c$.
  \item Compute $S$ (total sells) once.
  \item Per firm write $\min(b_c,\ S-s_c)$ --- ``wants vs.\ allowed''.
  \item Sum them; then $\min$ with $\min(B,S)$.
\end{enumerate}

\begin{say}
``Let me tabulate per company. Total sells is $S$. Each firm can buy at most what
it wants, $b_c$, and at most the sells that aren't its own, $S-s_c$ --- so its
reach is the smaller. I sum the reaches and cap by $\min(B,S)$ so I never exceed a
side. That handles the different-company rule exactly, and it reduces to naive
$\min$ when no firm dominates.''
\end{say}
